\chapter{Introduction et contexte}

\section{Contexte}

Une petite boutique qui gère ses produits et ses commandes à la main (fichiers Excel, échanges
téléphoniques) s'expose à des erreurs de stock, à une mauvaise traçabilité des commandes et à une
visibilité limitée auprès de ses clients. Le projet consiste à lui fournir une application web de vente en
ligne : un catalogue public, un espace client (panier, commandes, historique, profil) et un back-office
d'administration.

\section{Problématique}

Comment offrir à une boutique un outil web simple permettant de publier un catalogue organisé par
catégories, de recevoir des commandes en ligne et d'en suivre le traitement, tout en garantissant la
cohérence des données (stock, prix, historique) ?

\section{Objectifs}

\begin{enumerate}
  \item Permettre aux visiteurs de consulter le catalogue et de rechercher des produits.
  \item Permettre aux clients authentifiés de gérer un panier et de passer des commandes.
  \item Permettre à l'administrateur de gérer les catégories, les produits, les commandes et les utilisateurs.
  \item Concevoir le système avec UML avant son implémentation.
  \item Mettre en œuvre une architecture MVC en Jakarta EE (Servlets, JSP, JSTL) avec Hibernate pour la persistance.
  \item Versionner le travail avec Git et GitHub (tickets, branches, pull requests).
\end{enumerate}

\section{Périmètre}

Le périmètre est celui du cahier des charges (§1.4). Sont \textbf{hors périmètre} : le paiement en ligne réel
(le paiement est simulé, à la livraison), la livraison et le suivi transporteur, la gestion multi-boutiques,
l'application mobile et l'envoi d'e-mails transactionnels.

\section{Acteurs}

\begin{table}[H]
  \centering\small
  \caption{Acteurs du système}
  \begin{tabularx}{\linewidth}{P{2.4cm}LL}
    \ligneentete
    \entete{Acteur} & \entete{Description} & \entete{Accès} \\
    Visiteur & Internaute non authentifié & Catalogue, recherche, fiche produit, inscription, connexion \\
    \midrule
    Client & Utilisateur connecté avec le rôle CLIENT ; hérite des droits du Visiteur & Panier, commandes, historique, profil \\
    \midrule
    Administrateur & Utilisateur connecté avec le rôle ADMIN & Back-office : catégories, produits, stock, commandes, utilisateurs, tableau de bord \\
    \bottomrule
  \end{tabularx}
\end{table}
